\section{Further Work.}

The most immediate improvement of the graph-based RLA framework would be to adapt a test supermartingale to the non-linear margins of STV. 
This would allow for more efficient and flexible audits with arbitrary stopping times.

The shortcut of \S\ref{sec:shortcut} might also be extended to deal with cases where a strong candidate is seated in the middle of a batch elimination.

Another line to look into is the usage of audit graphs for partial audits of STV elections. 
Conspicuously absent from \S\ref{sec:results} are the 9-seat elections of Cambridge, MA: this is because their higher seat counts make the margin for the last seat systematically razor sharp.
In such a case audit graphs can still formulate an audit for the first 8 out of 9 seats.
The construction of a ``partially coherent'' graph might also allow for the study of per-candidate Margins of Victory.

More broadly, there is a vast canyon separating our understanding of how chaotic STV elections can be in the worst case and how well-behaved they tend to be in reality.
This paper puts a rope bridge over this canyon, but there is still a lot of work to be done in resolving the tension between theoretical brittleness and empirical stability of the Single Transferable Vote.
